\documentclass[11pt]{article}

\usepackage[utf8]{inputenc}
\usepackage[T1]{fontenc}
\usepackage{amsmath,amssymb,amsthm}
\usepackage{geometry}
\usepackage{xcolor}
\usepackage{listings}
\usepackage{hyperref}
\usepackage{enumitem}
\usepackage{booktabs}
\usepackage{url}

\geometry{margin=1in}
\hypersetup{colorlinks=true, linkcolor=blue!50!black,
            urlcolor=blue!50!black, citecolor=blue!50!black}

\lstdefinestyle{shell}{
    basicstyle=\ttfamily\footnotesize,
    backgroundcolor=\color{black!4},
    frame=single, framerule=0pt,
    breaklines=true, columns=fullflexible,
    keepspaces=true,
    showstringspaces=false,
    xleftmargin=0.5em, xrightmargin=0.5em,
}
\lstdefinestyle{cpp}{
    style=shell,
    language=C++,
    keywordstyle=\color{blue!50!black},
    commentstyle=\color{green!40!black},
    stringstyle=\color{purple},
}

\newcommand{\vE}{\mathbf{E}}
\newcommand{\vB}{\mathbf{B}}
\newcommand{\vD}{\mathbf{D}}
\newcommand{\vH}{\mathbf{H}}
\newcommand{\vJ}{\mathbf{J}}
\newcommand{\vx}{\mathbf{x}}
\newcommand{\hx}{\hat{x}}
\newcommand{\hy}{\hat{y}}
\newcommand{\hz}{\hat{z}}
\newcommand{\eps}{\varepsilon}

\title{\textbf{\texttt{Maxwell\_Kokkos}}\\
       Documentation for the Kokkos port of the SBP--SAT Maxwell solver}
\author{HPC course --- SMS 751/2255}
\date{April 2026}

\begin{document}

\maketitle
\tableofcontents
\bigskip

\begin{abstract}
\texttt{Maxwell\_Kokkos} is a three-dimensional finite-difference
solver for Maxwell's equations in a linear (possibly inhomogeneous)
dielectric, coupled to a hyperbolic constraint-damping system.  The
spatial discretization uses fourth-order summation-by-parts (SBP-4-2)
operators with one-sided closures at every physical face.  Boundary
conditions are imposed weakly via Simultaneous Approximation Term
(SAT) penalties on all six Cartesian faces.  The time integrator is
classical fourth-order Runge--Kutta.  The code is the Kokkos port of
the \texttt{Maxwell\_Penalty} reference: the physics and the
discretization are unchanged, but every compute kernel is expressed
as a \texttt{Kokkos::parallel\_for} or \texttt{parallel\_reduce} so
the same source builds for OpenMP on a CPU host and for CUDA on a
single GPU.  This document covers the physics, the numerical scheme,
the SBP-SAT penalty derivation, source models and their limitations,
the Kokkos integration, build and run instructions, tests, validation
status, and the appendix with explicit per-face SAT formulas.
\end{abstract}

\section{Overview}

\texttt{Maxwell\_Kokkos} is a Kokkos re-implementation of
\texttt{Maxwell\_Penalty}.  The physics, the
discretization, and the user-facing interface (TOML configuration,
HDF5 output, MPI Cartesian decomposition) are unchanged from the
parent code; the compute kernels have been ported to
\texttt{Kokkos::parallel\_for} and \texttt{Kokkos::parallel\_reduce}
so the same source builds for OpenMP on a CPU host and for CUDA on
a single GPU.  Companion documents:

\begin{itemize}[itemsep=2pt,topsep=2pt]
  \item \texttt{../Maxwell\_Penalty/CLAUDE.md} --- parent code
        notes (TOML schema, source types, lens demo, convergence
        rates, SBP--SAT sign convention warnings).
  \item \texttt{../Maxwell\_Penalty/doc/documentation.tex}
        --- the parent document this one is partially derived from.
        Many of the physics, characteristic, and limitations
        sections are imported verbatim from there.
  \item \texttt{../Maxwell\_Penalty/Develop\_Penalty.md} ---
        continuum derivation of the SBP--SAT penalty for the extended
        Maxwell system.
  \item \texttt{GPU\_Maxwell.md} --- the implementation plan that
        produced this code.
  \item \texttt{Report.md} --- focused report on the include and
        build-option design as actually shipped.
\end{itemize}

The remainder of this document is organised as follows.
Sections~\ref{sec:physics}--\ref{sec:material} cover the physics:
the extended-Maxwell system, the spatial discretization, the
SBP--SAT penalty, the available analytic sources, and the material
model.  Section~\ref{sec:limitations} catalogues the limitations of
characteristic-based absorbing boundaries and
Section~\ref{sec:pml} sketches the Perfectly Matched Layer extension
that would address them.
Section~\ref{sec:kokkos-arch} is the new value-add of this port:
how the Kokkos integration is laid out, the design choices that let
the existing SymPy-generated SAT macros work unchanged, and the
file map.
Sections~\ref{sec:build}--\ref{sec:running} cover building and
running the code, including the full TOML parameter schema and the
two canonical configurations.
Section~\ref{sec:tests} documents the test suite,
Section~\ref{sec:convergence} the convergence study (the strongest
correctness check), Section~\ref{sec:status} the overall validation
status, and Section~\ref{sec:known} the toolchain pitfalls
discovered when bringing the code up on an A100 host.
Appendix~\ref{app:sat-formulas} gives the explicit per-face SAT
formulas.

\section{The continuum system}
\label{sec:physics}

The code evolves the electromagnetic field $(\vE, \vB)$ in a
linear, possibly inhomogeneous, dielectric.  Let $\eps(\vx)$ be the
permittivity, $\mu(\vx)$ the permeability, and $\sigma(\vx)$ the
conductivity (Ohmic damping).  Define $\vD = \eps \vE$ and
$\vH = \vB / \mu$.  In Gaussian units, Maxwell's equations are
\begin{align}
  \partial_t \vD &= \nabla \times \vH - 4\pi \vJ, \label{eq:ampere}\\
  \partial_t \vB &= -\nabla \times \vE, \label{eq:faraday}\\
  \nabla \cdot \vD &= 4\pi \rho, \label{eq:gauss-e}\\
  \nabla \cdot \vB &= 0, \label{eq:gauss-b}
\end{align}
with Ohm's law $\vJ = \sigma \vE$.  Equations
\eqref{eq:ampere}--\eqref{eq:faraday} are the time-evolution
(\emph{hyperbolic}) equations; \eqref{eq:gauss-e}--\eqref{eq:gauss-b}
are the divergence constraints, preserved by the time evolution in
the continuum.

Numerically this preservation is not automatic.  Round-off and
truncation error in the discretized curl operators can seed small
constraint violations that, in a naive scheme, grow without bound.
The standard remedy --- following Munz et al.\ --- is to extend the
system with two auxiliary scalar fields $\Psi_D, \Psi_B$ that act as
Lagrange multipliers for the constraints:
\begin{align}
  \partial_t \vD - c^2 \nabla \Psi_D
                 &= \nabla \times \vH - 4\pi \vJ, \\
  \partial_t \vB - c^2 \nabla \Psi_B
                 &= -\nabla \times \vE, \\
  \partial_t \Psi_D &= \nabla \cdot \vD - 4\pi \rho - \kappa_D \Psi_D, \\
  \partial_t \Psi_B &= \nabla \cdot \vB - \kappa_B \Psi_B.
\end{align}
At $\Psi_D = \Psi_B = 0$ the original Maxwell system is recovered.
For $\kappa_{D,B} > 0$ the $\Psi$ fields satisfy damped wave
equations whose source is the constraint violation: any residual
violation radiates outward at speed $c$ and is exponentially damped
on timescale $1/\kappa$.  This absorbs round-off and truncation
noise, and absorbs the $O(1/(kw_0)^2)$ constraint residual of the
paraxial beam source (\S\ref{sec:gaussian-beam}).

The code evolves nine fields,
$U = (D_x, D_y, D_z, B_x, B_y, B_z, \Psi_D, \Psi_B, \rho)^\top$,
plus five auxiliaries $(1/\eps, 1/\mu, \sigma, c_D, c_B)$ where the
last two are diagnostic constraint values.  The charge density
$\rho$ is carried as an evolved variable for generality, driven by
$\partial_t \rho + \nabla \cdot \vJ = 0$; for the simulations
shipped with the code, $\rho$ and $\vJ$ are initially zero and
remain zero up to round-off.  Slot ordering is fixed in the
\texttt{evolved\_slot} and \texttt{auxilliary\_slot} enums in
\texttt{maxwell\_eqs.hpp}; the C-version's discipline of ``one
source of truth for variable order'' is preserved verbatim.

The material parameters $\eps, \mu, \sigma$ are stored in inverse
form $\mathtt{ieps} = 1/\eps$, $\mathtt{imu} = 1/\mu$,
$\mathtt{sigma} = \sigma$ as auxiliary gridfunctions.  They may
vary in space; the default beam-through-lens demo places an
elliptical-contour permittivity bump (\S\ref{sec:material}).

\subsection{The flux form}

The spatial derivatives in the extended system can be written as a
quasi-linear hyperbolic system
\begin{equation}
  \partial_t U + A_x \partial_x U + A_y \partial_y U + A_z \partial_z U
  = S(U, \vx),
  \label{eq:flux-form}
\end{equation}
where $S$ collects source terms (conductivity damping, constraint
relaxation, material-gradient couplings) and the three flux matrices
$A_x, A_y, A_z$ encode spatial propagation.  Their characteristic
decomposition is what the SAT penalty needs; we return to it in
\S\ref{sec:sat}.

\section{Spatial discretization}
\label{sec:spatial}

\subsection{Grid, decomposition, and indexing}

The domain is a rectangular parallelepiped
$[x_0, x_n]\times[y_0, y_n]\times[z_0, z_n]$, discretized on a
uniform Cartesian grid of global size $(N_x, N_y, N_z)$ with
spacings $dx, dy, dz$.  The grid is distributed across MPI ranks
via a Cartesian topology chosen automatically by greedy prime
factorization.  Each rank owns a contiguous block surrounded by
ghost zones of width $g_s$ (default 3).  For physical boundaries
the ghost zones are absent and the rank stores real data at every
index; for MPI-internal and periodic boundaries, the ghost zones
are synchronised via non-blocking point-to-point communication
between RK stages.

The local index into a 3D field $f$ at grid point $(i, j, k)$ is
$\mathtt{ijk} = i + j N_x + k N_x N_y$, with $i$ the fastest-varying
index.  Strides are $d_i = 1$, $d_j = N_x$, $d_k = N_x N_y$.  Every
finite-difference stencil takes a stride argument so it can operate
along any axis.  The Kokkos port keeps this offset arithmetic
identical: \texttt{Kokkos::View<double***, LayoutLeft>} of shape
$(N_x, N_y, N_z)$ has the same flat-index offset, so the existing
stencils and SAT macros work unchanged --- see
\S\ref{sec:kokkos-arch}.

\subsection{Summation-by-parts (SBP) operators}
\label{sec:sbp-operators}

The spatial derivatives are computed by summation-by-parts
finite differences of order $(4,2)$ --- fourth-order in the
interior, second-order one-sided closures at the first four rows
on each boundary.  A $(4,2)$ SBP operator $D$ approximating
$d/dx$ on a grid of $N$ points factorizes as
\begin{equation}
  D = H^{-1} Q, \qquad
  Q + Q^\top = B = \mathrm{diag}(-1, 0, \ldots, 0, +1),
\end{equation}
where $H$ is a positive-definite diagonal weight matrix (the
``norm'') and $Q$ is almost antisymmetric, the asymmetry captured
entirely by the boundary term $B$.  In the $(4,2)$ case
\begin{equation}
  H = h \cdot \mathrm{diag}(17/48,\, 59/48,\, 43/48,\, 49/48,\,
                            1, 1, \ldots).
\end{equation}
The key stability-enabling identity is
\begin{equation}
  \langle u, D v \rangle_H + \langle Du, v \rangle_H
    = u_N v_N - u_0 v_0,
\end{equation}
where
$\langle u, v \rangle_H = \sum_i H_{ii} u_i v_i$.  This is the
discrete analogue of integration by parts on a finite interval ---
the boundary terms appear exactly, without truncation error, and
this identity is what lets us prove a discrete energy estimate.

The interior stencil (rows $i = 4, \ldots, N-5$) is the standard
fourth-order centered difference
\[
  (D u)_i = \frac{u_{i-2} - 8 u_{i-1} + 8 u_{i+1} - u_{i+2}}{12 h}
  + O(h^4).
\]
The four near-boundary rows use one-sided Strand (1994) closures;
at the lower boundary,
\begin{align*}
  (D u)_0 &= \frac{-24 u_0 + 59 u_1 - 8 u_2 - 3 u_3}{17 h}, \\
  (D u)_1 &= \frac{-u_0 + u_2}{2 h}, \\
  (D u)_2 &= \frac{ 4 u_0 - 59 u_1 + 59 u_3 - 4 u_4}{86 h}, \\
  (D u)_3 &= \frac{ 3 u_0 - 59 u_2 + 64 u_4 - 8 u_5}{98 h},
\end{align*}
and symmetric closures at the upper boundary.  These are
implemented as \texttt{SBP42\_L0}--\texttt{SBP42\_L3} and the
upper-boundary counterparts \texttt{SBP42\_RN}--\texttt{SBP42\_RNm3}
in \texttt{derivatives.hpp}.  In the Kokkos port each stencil
function is annotated \texttt{KOKKOS\_INLINE\_FUNCTION} so it can
be called from device kernels.

\subsection{Optional Kreiss--Oliger dissipation}

When \texttt{[solver].use\_dissipation = true} the RHS is augmented
by a 5th-order Kreiss--Oliger dissipation operator
\[
  \dot u_i \mathrel{+}=
    -\frac{\sigma_{\text{KO}}}{\Delta x}
    \cdot \frac{1}{64}\bigl(
       u_{i-3} - 6 u_{i-2} + 15 u_{i-1}
       - 20 u_i
       + 15 u_{i+1} - 6 u_{i+2} + u_{i+3}\bigr),
\]
applied separately along $x$, $y$, $z$.  The operator is skipped
on the outer three rows of each axis.  Useful for damping
grid-scale noise from oblique partial reflections; see
\S\ref{sec:oblique-reflection}.

\section{Time integration}
\label{sec:time}

Time integration is classical explicit fourth-order Runge--Kutta:
\begin{align*}
  k_1 &= f(t_n,\, u_n),                  & u_n^{(1)} &= u_n + \tfrac{\Delta t}{2} k_1,\\
  k_2 &= f(t_n + \tfrac{\Delta t}{2},\, u_n^{(1)}), &
                     u_n^{(2)} &= u_n + \tfrac{\Delta t}{2} k_2,\\
  k_3 &= f(t_n + \tfrac{\Delta t}{2},\, u_n^{(2)}), &
                     u_n^{(3)} &= u_n + \Delta t\, k_3,\\
  k_4 &= f(t_n + \Delta t,\, u_n^{(3)}), &
                     u_{n+1} &= u_n + \tfrac{\Delta t}{6}\,
                        (k_1 + 2 k_2 + 2 k_3 + k_4).
\end{align*}
The timestep is chosen by the CFL condition
\(\Delta t = c \cdot \min(\Delta x, \Delta y, \Delta z)\)
with the safety factor $c$ taken from the TOML
(\texttt{[solver].cfl\_factor}).  For the pure SBP interior (no
SAT), stability permits $c \lesssim 1$.  With SAT penalties on all
six faces, the effective spectral radius increases near corners
where multiple penalties stack, and the usable CFL drops to
$c \approx 0.2$ in practice (\S\ref{sec:stability-tau}).

In the Kokkos port the C version's pointer-rebinding trick (where
each \texttt{vars[v]->dot} pointer was reseated between RK stages
to alias \texttt{K1..K4}) is replaced by an explicit integer
\texttt{kidx} parameter to the time-derivative function.
\texttt{DECLARE\_EVOLVED\_VARS} uses \texttt{kidx} to pick the
correct \texttt{K[kidx].data()} pointer.  This keeps the kernel
body identical to the C version while removing a non-obvious
dataflow.

\section{Boundary conditions: the SBP-SAT penalty method}
\label{sec:sat}

\subsection{Why weak imposition}

The naïve approach to boundary conditions --- ``overwrite the
evolution at boundary nodes with the prescribed boundary value''
--- is called \emph{strong imposition}.  It is simple and often
works, but it breaks the SBP energy estimate: a boundary row whose
equation is replaced loses the clean integration-by-parts
structure, and stability must be argued separately, usually by
numerical experiment.

The Simultaneous Approximation Term (SAT) method imposes boundary
conditions \emph{weakly}, as an additive penalty.  The PDE is
retained at every row, including the boundary row, and a correction
is \emph{added} that pulls the solution toward the prescribed
value.  The SAT term is tuned so the full semi-discrete system
admits a discrete energy estimate: summing the equations against
$U^\top H$ and using the SBP identity produces boundary terms that
are exactly cancelled by the SAT penalty.  This gives provable
stability for any $\tau \geq 1/2$
(\S\ref{sec:stability-tau}), independent of the particular PDE
coefficients or grid spacing.

\begin{center}
\small
\begin{tabular}{lll}
\toprule
Property                       & Strong imposition            & SBP-SAT (weak)      \\
\midrule
Boundary equation              & Replaced entirely            & Interior RHS + penalty \\
Energy estimate                & No formal proof              & Yes, provable for $\tau \geq 1/2$ \\
Stability w/ non-uniform $\eps, \mu$ & No guarantee           & Provable \\
Convergence order              & Same as interior closures    & Same as interior closures \\
Penalty scaling                & N/A                          & $\tau \cdot H^{-1}_{00}/h$ \\
\bottomrule
\end{tabular}
\end{center}

Both methods give the same formal convergence order --- the
boundary truncation error is set by the SBP one-sided closure, not
by how the boundary condition is imposed.  The pedagogical reason
to prefer SBP-SAT is the discrete energy estimate, which extends
to variable coefficients, multi-dimensional grids, and
characteristic couplings.

\subsection{Characteristic decomposition}
\label{sec:characteristics}

The penalty acts on the \emph{incoming} characteristic at each
face.  For a hyperbolic system $\partial_t U + A \partial_x U = 0$
with flux matrix $A$, diagonalize
\[
  A = V \Lambda V^{-1},
  \qquad \Lambda = \mathrm{diag}(\lambda_1, \ldots, \lambda_n),
\]
and split into positive and negative eigenvalue parts:
\[
  A^+ = V \Lambda^+ V^{-1}, \quad
  A^- = V \Lambda^- V^{-1}, \quad
  \Lambda^\pm = \mathrm{diag}(\max(\pm \lambda_i, 0)).
\]
The characteristic variables $W = V^{-1} U$ satisfy
$(\partial_t + \lambda_i \partial_x) W_i = 0$; modes with
$\lambda_i > 0$ propagate in $+x$ (incoming at the lower face),
modes with $\lambda_i < 0$ propagate in $-x$ (incoming at the
upper face).

At the lower-$x$ face, we prescribe the incoming characteristics
(the $+\lambda$ modes); the $-\lambda$ modes are outflow and are
determined by the interior dynamics.  At the upper-$x$ face the
roles are reversed.

\subsection{The penalty form}

At the lower-$x$ boundary (grid index $i = 0$), the SAT correction
is
\begin{equation}
  \dot U_0 \;\mathrel{-}=\;
   \tau \cdot \frac{(H^{-1})_{00}}{h} \cdot A^+ \cdot (U_0 - g),
  \label{eq:sat-lower}
\end{equation}
where $g$ is the prescribed incoming state and $\tau \geq 1/2$ is
the penalty strength.  The code uses $\tau = 1$ by default;
$(H^{-1})_{00} = 48/17$ for the $(4,2)$ SBP norm.  The stencil
operates only on the single boundary row --- rows
$i = 1, 2, 3$ receive the SBP one-sided derivative
(\S\ref{sec:sbp-operators}) but \emph{no} additional SAT term;
they contribute to the energy argument through the norm integral.
The analogous correction at the upper boundary uses $A^-$ (or
equivalently the positive-definite projection $|A^-|$) and the
incoming data for that face:
\begin{equation}
  \dot U_{N-1} \;\mathrel{-}=\;
   \tau \cdot \frac{(H^{-1})_{N-1,N-1}}{h}
   \cdot |A^-| \cdot (U_{N-1} - g).
  \label{eq:sat-upper}
\end{equation}
By symmetry $(H^{-1})_{N-1,N-1} = (H^{-1})_{00} = 48/17$.

\subsection{Application to Maxwell's equations}

For the extended Maxwell system, the $z$-axis flux matrix
decomposes into four independent $2\times 2$ blocks plus one
trivial $1\times 1$ block:

\begin{center}
\small
\begin{tabular}{llll}
\toprule
Block        & Variables & Eigenvalues & Incoming char.\ at lower $z$ ($w^+$) \\
\midrule
EM-1         & $(B_x, D_y)$     & $\pm c$ & $B_x - PP \cdot D_y$ \\
EM-2         & $(B_y, D_x)$     & $\pm c$ & $B_y + PP \cdot D_x$ \\
Div-$B$      & $(B_z, \Psi_B)$  & $\pm 1$ & $B_z - \Psi_B$ \\
Div-$D$      & $(D_z, \Psi_D)$  & $\pm 1$ & $D_z - \Psi_D$ \\
Charge       & $\rho$           & $0$     & (none) \\
\bottomrule
\end{tabular}
\end{center}

Here $c = 1/\sqrt{\eps\mu}$ is the local light speed
($\sqrt{\mathtt{ieps}\cdot\mathtt{imu}}$ in the code's
normalisation) and $PP = \sqrt{\eps/\mu}^{-1}
= \sqrt{\mathtt{ieps}/\mathtt{imu}}$ is the impedance ratio.  The
divergence blocks have eigenvalues $\pm 1$ because constraint
violations propagate at the normalised speed set by the extended
system, independent of material.  The $w^+$ column lists, for
each block, the left-eigenvector combination that is the incoming
characteristic at the lower-$z$ face; the upper-$z$ face injects
the opposite combination ($w^-$) with the signs of the
cross-coupling terms reversed.

For the EM-1 block $(B_x, D_y)$, the flux matrix is
\[
  A_z^{(\mathrm{EM1})}
  = \begin{pmatrix} 0 & -\mathtt{ieps} \\
                    -\mathtt{imu} & 0 \end{pmatrix}.
\]
The right eigenvector for $\lambda = +c$ is
$v^+ = (1, -1/PP)^\top$; the left eigenvector
($L^+ A = c L^+$) is $L^+ = (1, -PP)^\top$; the positive
projection is
\[
  A^+ = \frac{c}{L^+ \cdot v^+}\, v^+ (L^+)^\top
      = \frac{c}{2}
        \begin{pmatrix} 1 & -PP \\ -1/PP & 1 \end{pmatrix}.
\]
Applying to the deviation
$\delta U = (B_x - g_{B_x}, D_y - g_{D_y})$:
\begin{align*}
  (A^+ \delta U)_{B_x}
    &= \tfrac{c}{2}(\delta B_x - PP \cdot \delta D_y), \\
  (A^+ \delta U)_{D_y}
    &= \tfrac{c}{2}(-\delta B_x / PP + \delta D_y).
\end{align*}
The incoming characteristic for this block is therefore
$w^+ = L^+ \cdot U = B_x - PP \cdot D_y$.

\paragraph{Sanity check on a known solution.}
For a $y$-polarised plane wave travelling in $+z$ (incoming at the
lower-$z$ boundary), the fields satisfy
$D_y = \psi$, $B_x = -\psi$ (the sign makes $\vE \times \vB$ point
in $+\hz$), so
\[
  w^+ = -\psi - PP\cdot\psi = -(1+PP)\psi \neq 0
\]
as expected for an incoming mode.  A sign error that flipped any
coupling term ($+PP$ instead of $-PP$, say) would give $w^+ = 0$
on this same physical solution, and the SAT would silently fail to
inject the wave.  Earlier hand-written sign conventions for this
system are known to contain errors of exactly this form; the
SymPy-emitted macros in \texttt{simple\_maxwell.h} are the
authoritative version, derived from scratch by SymPy
diagonalization of $M_a$ and verified against the traveling-wave
check above.

The implementation lives in \texttt{simple\_maxwell.h} as six
preprocessor macros: \texttt{APPLY\_SAT\_LOWER\_\{X,Y,Z\}} and
\texttt{APPLY\_SAT\_UPPER\_\{X,Y,Z\}}.  Only the lower-$z$ face
has nonzero $g$ in the plane-wave and Gaussian-beam configurations
(the incoming source field, \S\ref{sec:sources}); the other five
faces use $g = 0$, which corresponds to a perfectly absorbing
boundary for waves leaving normally.  See
\S\ref{sec:oblique-reflection} for the angular caveat.

\subsection{Edges and corners: additivity}

At an edge (two physical faces meet) and at a corner (three faces
meet), the SAT penalties are \emph{additive}: each face
contributes its own penalty term independently.  No cross-term or
corner-specific correction is needed.  This additivity is a direct
consequence of the tensor-product structure: the 3D discrete
energy norm $H = H_x \otimes H_y \otimes H_z$ decomposes the
boundary contributions into independent surface integrals over
each face, each cancelled by its own SAT term.

Concretely, at the corner $(i,j,k) = (0,0,0)$ with all three axes
physical, the RHS receives
\begin{equation}
  \dot U_{0,0,0} \;\mathrel{-}=\; \tau \left[
    \frac{H^{-1}_0}{dx} A_x^+ \delta U
  + \frac{H^{-1}_0}{dy} A_y^+ \delta U
  + \frac{H^{-1}_0}{dz} A_z^+ \delta U
  \right],
  \label{eq:corner-sat}
\end{equation}
where the three projections act on the same state $\delta U$ but
with axis-specific flux matrices.  This is the only corner-aware
term in the code --- the SBP derivatives at the corner use
one-sided closures in each direction independently, with no other
correction.

\subsection{Stability, $\tau$, and CFL at corners}
\label{sec:stability-tau}

The standard SBP-SAT stability theorem (Svärd \& Nordström)
requires $\tau \geq 1/2$ for a single-face penalty on a scalar
hyperbolic equation.  For a system with $n$ characteristic speeds
on a tensor-product multidimensional grid, the same bound applies
per face.  The code uses $\tau = 1$ as a comfortable margin.

At corners, three SAT terms with per-face scaling
$\tau (H^{-1}_0)/h$ stack additively on the same boundary row.
The effective spectral radius of the discrete operator grows at
the corner, and the usable CFL number drops.  Empirically, the
six-face-physical scheme is stable at $\mathrm{cfl} \approx 0.2$
with classical RK4 --- tighter than $\mathrm{cfl} \approx 0.5$
safe for an interior-only (or two-face) scheme.  Raising $\tau$
beyond 1 tightens the CFL further; reducing $\tau$ toward $1/2$
relaxes it at the cost of stability margin.

\section{Source models}
\label{sec:sources}

The lower-$z$ face (or, for the convergence-test source, every
physical face) has nonzero boundary data $g$ corresponding to a
specified incoming electromagnetic field.  Three sources are
implemented:

\begin{itemize}[itemsep=2pt,topsep=2pt]
  \item \texttt{plane\_wave}: a bump-enveloped plane wave
        travelling in $+\hz$ (\S\ref{sec:plane-wave}).  Pair with
        periodic $x, y$ for the formal convergence reference
        (Option B).
  \item \texttt{gaussian\_beam}: a paraxial collimated beam
        injected at the lower-$z$ face (\S\ref{sec:gaussian-beam}).
        Used with all faces non-periodic for the headline
        beam-through-lens demonstration (Option A).
  \item \texttt{te\_waveguide\_mode}: an exact vacuum Maxwell TE
        mode of a rectangular unit-cube waveguide
        (\S\ref{sec:te-mode}), used as SAT data on every physical
        face for the four-configuration convergence study
        (\S\ref{sec:convergence}).
\end{itemize}
Selection is via \texttt{[source].type} in the TOML.  Per-source
parameters live in dedicated sub-tables; the driver parses all
three so switching sources is a one-line edit.

\subsection{Plane wave source}
\label{sec:plane-wave}

A transverse electromagnetic wave travelling in $+\hz$, linearly
polarised as
$\vE = (E_0^x \hx + E_0^y \hy) f(t-z) \cos(k(z-t))$, modulated by
a $C^2$ smooth bump envelope $f$ that turns on over an interval
$[a, b]$.  This satisfies the vacuum Maxwell equations exactly
(independent of $k$), so in a configuration with $x, y$ periodic
and $z$ non-periodic this is an \emph{exact} reference for formal
convergence testing of the SBP-SAT scheme.

The plane wave has no transverse structure: the field at $x = 0$
and $x = L_x$ is identical.  Combining this source with
non-periodic $x, y$ would put a full-amplitude wavefield at the
transverse walls, where the $g = 0$ SAT pulls toward zero and
excites a numerical instability.  The plane wave is not usable
with six-face SAT (\S\ref{sec:limitations}).

\subsection{Paraxial Gaussian beam}
\label{sec:gaussian-beam}

Models a collimated laser beam entering at the lower-$z$ face,
optionally converging toward a waist at $z = z_w$.  The paraxial
scalar amplitude is
\begin{equation}
  \psi(\vx, t) = E_0 \cdot \frac{w_0}{w(z)}
    \exp\!\left(-\frac{r^2}{w(z)^2}\right)
    \cos\!\left(kz - \omega t + \frac{k r^2}{2 R(z)} - \varphi_G(z)\right),
  \label{eq:gauss-amp}
\end{equation}
with $r^2 = x^2 + y^2$, $\omega = kc$, and
\begin{align*}
  z_R &= \tfrac{1}{2} k w_0^2
            \quad\text{(Rayleigh range)} \\
  w(z)       &= w_0 \sqrt{1 + ((z-z_w)/z_R)^2} \\
  1/R(z)     &= (z-z_w) / ((z-z_w)^2 + z_R^2) \\
  \varphi_G(z)& = \arctan((z-z_w)/z_R)
            \quad\text{(Gouy phase)}.
\end{align*}
The implementation in \texttt{analytic\_solutions.hpp} uses an
electric Hertz-vector construction
$\boldsymbol{\Pi} = \hat{x}\,\mathrm{Re}[\Phi_c]$ with
$\Phi_c = (E_0/k^2)\, A(r, z)\, e^{i\gamma}$ and sets
$\vD = \nabla\times\nabla\times\boldsymbol{\Pi}$,
$\vB = \partial_t(\nabla\times\boldsymbol{\Pi})$.  This makes
$\nabla\cdot\vD = \nabla\cdot\vB = 0$ \emph{exactly} by
$\nabla\cdot(\nabla\times\cdot) \equiv 0$, so the constraint-damping
fields stay at numerical-noise level even though $\Phi$ only solves
the wave equation to $O((kw_0)^{-2})$.

\paragraph{Convergence vs.\ divergence.}
The sign of $1/R(z)$ determines the wavefront curvature.  For
$z < z_w$, $R < 0$: the wavefronts are concave toward the beam
axis and the beam converges.  For $z > z_w$, $R > 0$ and the beam
diverges.  Placing $z_w$ in the middle of the domain produces a
beam that converges through the first half, reaches minimum spot
at the middle, then diverges.

\paragraph{Collimation.}
A beam is ``nearly collimated'' (laser-like) over a propagation
distance $L$ when $z_R \gg L$.  The spot size varies as
$w(L)/w_0 = \sqrt{1 + (L/z_R)^2}$, so $z_R = 10 L$ gives 0.5\,\%
broadening, and $z_R = 3 L$ gives 5\,\%.  Fixed domain $L$ then
requires large $k w_0^2$, which costs transverse resolution
(\S\ref{sec:limitations}).

\paragraph{Paraxial residual.}
\eqref{eq:gauss-amp} satisfies the full Maxwell equations only to
$O(1/(kw_0)^2)$.  The residual manifests as small but nonzero
$\nabla\cdot\vE$ and $\nabla\cdot\vB$ when the beam is injected;
the constraint-damping fields $\Psi_D, \Psi_B$ absorb and radiate
it on timescale $1/\kappa_{D,B}$.  $kw_0 \gtrsim 5$ keeps the
residual below $4\%$, well within the constraint-damping system's
domain of validity.

\subsection{Smooth turn-on in time}

For \texttt{plane\_wave}, the envelope $f(t-z)$ vanishes at $t=0$
for all $z\geq 0$, so the initial state is identically zero and
the wave enters gradually through the SAT boundary.

For \texttt{gaussian\_beam}, the analytic formula
\eqref{eq:gauss-amp} does not vanish at $t = 0$.  Filling the
interior with it would seed the domain with a state that is only
a paraxial approximation; empirically this excites a numerical
instability that grows exponentially regardless of CFL.  To avoid
it, the beam source is multiplied by a smooth $C^2$ turn-on
$f(t; t_a, t_b)$ that rises from 0 at $t \leq t_a$ to 1 at
$t \geq t_b$.  The interval is exposed as
\texttt{[source.gaussian\_beam].ramp\_a}/\texttt{ramp\_b};
default $[0,1]$.  The initial state is therefore identically
zero, and the beam grows in through the boundary SAT as the
envelope rises.  Lengthening the ramp reduces paraxial-residual
excitation during turn-on at the cost of delaying the onset of
interesting physics.

The \texttt{te\_waveguide\_mode} source is an exact solution and
needs no turn-on: it is non-zero at $t = 0$ and the initial state
on the full grid is consistent with the boundary data at every
face.  The convergence test relies on this consistency
(\S\ref{sec:convergence}).

\subsection{TE waveguide mode}
\label{sec:te-mode}

The TE($l, m, n$) mode of a unit-cube rectangular waveguide
propagating in $+z$, with dispersion
$\omega^2 = \pi^2(l^2 + m^2 + n^2)$:
\begin{align*}
  D_x &= \frac{\omega}{l\pi}\, \sin(m\pi y)\, \cos(l\pi x)\, \cos(n\pi z - \omega t), \\
  D_y &= -\frac{\omega}{m\pi}\, \sin(l\pi x)\, \cos(m\pi y)\, \cos(n\pi z - \omega t), \\
  D_z &= 0, \\
  B_x &= \frac{n}{m}\, \sin(l\pi x)\, \cos(m\pi y)\, \cos(n\pi z - \omega t), \\
  B_y &= \frac{n}{l}\, \sin(m\pi y)\, \cos(l\pi x)\, \cos(n\pi z - \omega t), \\
  B_z &= -\frac{l^2+m^2}{lm}\, \sin(n\pi z - \omega t)\, \cos(l\pi x)\, \cos(m\pi y),
\end{align*}
with $\Psi_D = \Psi_B = \rho = 0$.  The mode numbers
$(l, m, n)$ must each be $\geq 1$ (the parameter parser enforces
this; division by $l$, $m$ in the formula above also requires it).
For even $(l, m, n)$ the solution is also periodic on the unit
cube, so the same reference works for both periodic and
fully-physical boundary configurations in the convergence study.

\section{Material model}
\label{sec:material}

\paragraph{Background.}
A uniform $(\eps, \mu, \sigma)$ from
\texttt{[material.background]}.  The auxiliary fields
$1/\eps, 1/\mu, \sigma$ live in slots
\texttt{IEPS\_SLOT, IMU\_SLOT, SIGMA\_SLOT} and are constant in
time.

\paragraph{Optional elliptical lens.}
If \texttt{[material].epsilon\_type = "elliptical"} an
elliptical-Gaussian permittivity profile is layered on top:
\begin{equation}
  \eps(x, z) = (\eps_{\max} - \eps_{\text{bg}})\,
    \exp\!\left[-s^2 \left(\frac{(x-x_0)^2}{a^2}
                          + \frac{(z-z_0)^2}{b^2}\right)^2\right]
    + \eps_{\text{bg}},
  \label{eq:lens}
\end{equation}
with $\eps_{\text{bg}}$ the \texttt{epsilon} from
\texttt{[material.background]}.  Set once at
\texttt{set\_initial\_data} time and held constant thereafter.
Used in the Option A demo.

\section{Limitations}
\label{sec:limitations}

\subsection{Transverse-face reflections}

The $g = 0$ SAT on the $\pm x, \pm y$ faces is correct only
insofar as the physical solution has negligible amplitude there.
For a Gaussian beam with $w_0 \ll L_x, L_y$ the wall tail decays
as $\exp(-(L_{x,y}/2 w)^2)$ and reflections sit below roundoff for
$L_{x,y} \gtrsim 10\, w_0$.  For a plane wave (full amplitude at
the walls) this boundary treatment is unsuitable --- spurious
reflection cascades into instability.

\subsection{Paraxial approximation}

The Gaussian-beam source is a paraxial approximation to Maxwell's
equations.  The residual constraint violation is $O(1/(kw_0)^2)$
and is absorbed by the constraint-damping fields.  This is
acceptable as a physics source but means the beam is not an exact
solution; the $L^2$ error of the numerical evolution against the
paraxial reference is therefore \emph{not} a convergence metric.
Use the plane-wave (\S\ref{sec:plane-wave}) or
\texttt{te\_waveguide\_mode} (\S\ref{sec:te-mode}) source for
convergence testing.

\subsection{Resolution / containment / paraxial trade-off}

Three constraints on the beam parameters:
\begin{itemize}[itemsep=1pt,topsep=2pt]
  \item Grid resolution: $k\cdot dx \lesssim 0.6$, i.e., roughly
        10 points per wavelength.
  \item Paraxial validity: $k w_0 \gtrsim 5$ (residual $\leq 4\%$).
  \item Transverse containment: $L_x/w_0 \gtrsim 5$
        (preferably $\geq 10$).
\end{itemize}
For a laser-like (collimated) beam, the Rayleigh range
$z_R = k w_0^2 / 2$ must additionally be large compared with the
propagation length $L_z$.  All four together favour large $k$ and
large $w_0$, raising the required $N_x$ steeply.

\subsection{CFL at corners}

With all six faces treated by SAT, the scheme is stable at
$\mathrm{cfl} \approx 0.2$ with classical RK4, tighter than the
$\mathrm{cfl} \approx 0.5$ safe for the periodic-$xy$
interior-only scheme (\S\ref{sec:stability-tau}).  Exponentially
growing $L^2$ error usually means the CFL is too large, not a
physics failure.

\subsection{Plane wave with six-face SAT}

The plane-wave source combined with all six SAT faces produces a
documented numerical instability: the transverse walls see
full-amplitude tangential field, and the $g = 0$ SAT attempts to
force it to zero.  This is not a bug; it is the documented
failure mode of the SAT boundary for a physical configuration
whose field is not transversely contained.  Use periodic $x, y$
for plane-wave runs.

\subsection{Angular (oblique) reflections from every SAT face}
\label{sec:oblique-reflection}

The most important limitation of the current boundary treatment is
shared by all characteristic-based absorbing boundaries: the SAT
penalty is exact for waves moving \emph{normal} to the boundary,
but only partial for waves moving at oblique angles.  This
affects every face, not just the transverse ones:
\begin{itemize}[itemsep=2pt,topsep=2pt]
  \item \textbf{Lower-$z$ face}: the SAT's $+c$ projection
        correctly leaves a head-on outgoing wave (the $-c$ mode)
        untouched.  A wave at angle $\theta$ to the $z$-axis has
        its energy split between the $+c$ and $-c$
        $z$-characteristics in proportions depending on
        $\sin\theta$; the SAT treats the $+c$ part as incoming and
        pulls it toward the prescribed beam data, producing a
        spurious reflection whose amplitude grows with $\theta$.
  \item \textbf{Upper-$z$ face}: same mechanism, with $g = 0$.  A
        wave refracted by the lens and travelling mostly in $+z$
        but at an angle is partially reflected back into the
        domain.
  \item \textbf{$\pm x, \pm y$ faces}: same mechanism again.  The
        beam reaches these walls only as a Gaussian tail and as
        lens-scattered components, so the contribution is small
        for well-contained beams but non-negligible whenever the
        lens refracts energy sideways.
\end{itemize}

\paragraph{Concrete failure mode: beam on a high-$\eps$ lens.}
A dielectric step of permittivity $\eps_{\max}$ reflects a Fresnel
fraction $R = ((\sqrt{\eps_{\max}} - 1)/(\sqrt{\eps_{\max}} + 1))^2$
of incident intensity at normal incidence.  For $\eps_{\max} = 4$,
$R \approx 0.11$ in intensity, $\sim 1/3$ in amplitude --- a
substantial portion of the beam bounces off the lens and travels
back toward $-z$.  The pure-$-z$ component exits cleanly; angular
side-lobes generated by the elliptical edge of the lens partially
reflect off the SAT faces and re-enter.  After a few bounces the
domain fills with low-amplitude oblique modes that interfere with
the incident beam every $\lambda/2 = \pi/k$.

\paragraph{Mitigations within the current code.}
\begin{itemize}[itemsep=2pt,topsep=2pt]
  \item Shorten the simulation time so no wave completes a round
        trip.
  \item Weaken the lens ($\eps_{\max}$ closer to 1).
  \item Widen the domain so the transverse walls stay far from
        the beam after lens refraction.
  \item Enable Kreiss-Oliger dissipation
        (\texttt{use\_dissipation = true},
        \texttt{diss\_coeff} $\sim 0.02$--$0.1$) to damp grid-scale
        noise from partial reflections.  This does not remove the
        angular reflections themselves but suppresses their
        high-frequency signature.
\end{itemize}
None of these is a true fix; for angle-independent absorption the
code would need a Perfectly Matched Layer (\S\ref{sec:pml}).

\section{Perfectly Matched Layers: a future extension}
\label{sec:pml}

A Perfectly Matched Layer (PML) is a finite-thickness absorbing
buffer surrounding the physical domain, designed to absorb
outgoing waves \emph{at any angle and any frequency} with zero
reflection at the interface (in the continuum).  PML is the
standard tool in production electromagnetic codes for open-domain
problems.  \texttt{Maxwell\_Kokkos} does \emph{not} include a PML;
this section sketches what one would add.

\subsection{The mathematical mechanism: complex coordinate stretching}

For Maxwell in, say, an $x$-face PML, formally replace
$x$-derivatives by stretched derivatives
\begin{equation}
  \frac{\partial}{\partial x} \;\longrightarrow\;
  \frac{1}{s_x(x)} \frac{\partial}{\partial x},
  \qquad
  s_x(x) = 1 + \frac{\sigma_x(x)}{i\omega},
  \label{eq:pml-stretch}
\end{equation}
with $\sigma_x(x) \geq 0$ the PML damping profile, zero in the
physical domain and ramping up through the PML.  Substituting a
plane wave $e^{i(kx - \omega t)}$ in the PML region gives
\[
  e^{i(kx - \omega t)}\,
  \exp\!\left(-\int_{x_0}^x \frac{\sigma_x(x')}{c}\, dx'\right),
\]
the same wave with an exponential envelope.  Crucially, the
wavenumber $k$ and phase velocity are unchanged, so the wave does
not refract at the PML interface --- the match is perfect for
\emph{all} angles and \emph{all} frequencies.

\subsection{From frequency domain to time domain}

The complex factor $1/s_x$ is trivial in frequency but requires
care in time: dividing by $(i\omega + \sigma_x)$ is a convolution
with an exponential kernel.  Two standard implementations:
\begin{itemize}[itemsep=2pt,topsep=2pt]
  \item \textbf{Bérenger split-field PML (1994)}: each field
        component is split into pieces labelled by the contributing
        derivative; each piece evolves under its own damping ODE.
        Simple but doubles the field count in the PML.
  \item \textbf{Convolutional PML / CPML (Roden \& Gedney 2000)}:
        same complex stretch, but the convolution is implemented
        through one auxiliary field per stretched derivative,
        evolved by a first-order ODE.  Far fewer extra fields,
        more robust at grazing incidence, gracefully handles
        inhomogeneous media (like our lens).  CPML is the modern
        standard.
\end{itemize}
In CPML, each stretched derivative $(1/s_x)\partial_x E$ is
replaced by $\partial_x E + \Psi^x$ where the auxiliary $\Psi^x$
evolves as
\[
  \frac{d\Psi^x}{dt} = -\alpha_x \Psi^x - \sigma_x\, \partial_x E.
\]
Outside the PML, $\sigma_x = \alpha_x = 0$ and $\Psi^x \equiv 0$.

\subsection{The damping profile and what a PML would add}

$\sigma(x)$ starts at zero at the interface and rises smoothly to
$\sigma_{\max}$ at the outer face:
\(
  \sigma(x) = \sigma_{\max}\,(\xi/L_{\text{PML}})^m,
\)
$m = 3$ typical, with
$\sigma_{\max}$ tuned so the round-trip reflection
$R = \exp(-2\sigma_{\max} L_{\text{PML}}/((m+1)c))$ sits at the
target ($R \sim 10^{-5}$).  Adding CPML to
\texttt{Maxwell\_Kokkos} would require, roughly:
\begin{enumerate}[itemsep=1pt,topsep=2pt]
  \item Pad the grid by $N_{\text{PML}} \approx 10$ cells per
        physical face.
  \item Add six auxiliary $\Psi$ fields per spatial derivative on
        the PML region.
  \item Replace \texttt{SIMPLE\_MAXWELL\_INTERIOR\_DOT} with a
        PML-aware cousin in the layer cells (a per-point helper
        function rather than a fixed-coefficient macro).
  \item Apply matching damping profiles to the constraint-damping
        $\Psi_D, \Psi_B$ blocks.
  \item Use a trivial zero-Dirichlet (or zero-SAT) boundary at the
        outermost cell of the PML.
\end{enumerate}
A few hundred lines of code and a week of careful convergence
testing.  Comparison summary:

\begin{center}
\small
\begin{tabular}{lll}
\toprule
                   & SBP-SAT (current)                & PML (future) \\
\midrule
Mechanism          & Characteristic projection         & Complex coordinate stretching \\
Normal incidence   & Exact for $\tau \geq 1/2$         & Near-exact, $R \sim 10^{-5}$ \\
Grazing incidence  & Strongly reflective               & Near-exact \\
Extra grid cells   & None                              & $\sim 10$ per physical face \\
Extra fields       & None                              & One $\Psi$ per stretched derivative \\
Energy estimate    & Yes (provable)                    & Indirect (round-trip bound) \\
Handles lens refl. & Head-on only                      & All angles \\
Code complexity    & Low                               & Medium--high \\
\bottomrule
\end{tabular}
\end{center}

The choice between SAT and PML is a design decision, not a
correctness one.  For a teaching code demonstrating SBP-SAT, SAT
alone is appropriate.  For production scattering simulations, PML
is the right tool.

\section{Software architecture: Kokkos integration}
\label{sec:kokkos-arch}

This section is the unique value-added of the Kokkos port.
Sections~\ref{sec:physics}--\ref{sec:pml} apply equally to the
parent C version; the rest of this document is specific to
\texttt{Maxwell\_Kokkos}.

\subsection{Driving design choice}

A \texttt{Kokkos::View<double***, LayoutLeft>} sized
$(N_x, N_y, N_z)$ has the same flat-index offset
$\mathtt{ijk} = i + j N_x + k N_x N_y$ as the C version's
\texttt{double*} arrays.  Every kernel extracts a raw
\texttt{double*} via \texttt{view.data()} \emph{outside} the
parallel region:

\begin{lstlisting}[style=cpp]
double *Bx    = gfs->evol[BX_SLOT].state.data();
double *dotBx = gfs->evol[BX_SLOT].K[kidx].data();
// ...

Kokkos::parallel_for("interior", policy,
    KOKKOS_LAMBDA(int64_t i, int64_t j, int64_t k) {
        const int64_t ijk = i + j*nx + k*nx*ny;
        SIMPLE_MAXWELL_INTERIOR_DOT;   // uses Bx[ijk], dotBx[ijk], ...
    });
\end{lstlisting}

The lambda captures the pointers \emph{by value} (\texttt{[=]}); on
the OpenMP backend they point to host memory, on CUDA to device
memory.  Either way, \texttt{Bx[ijk]} resolves through the captured
pointer to the right address.

\paragraph{The single payoff.}
The SymPy-generated \texttt{simple\_maxwell.h} macros and the
SBP-4-2 stencils in \texttt{derivatives.hpp} need no source changes
other than annotating each \texttt{static inline} stencil with
\texttt{KOKKOS\_INLINE\_FUNCTION}.  The pointer-aliasing strategy
keeps the access pattern \texttt{Bx[ijk]} working unchanged --- so
\texttt{generate\_ccode.py} did not need updating, and
\texttt{make check-generated} (the SymPy-vs-source consistency
check from the parent) still passes byte-for-byte.

\subsection{Three-region kernel partition}

The time-derivative function \texttt{maxwell\_eq\_time\_deriv}
issues four \texttt{parallel\_for} launches, mirroring the C
version's three-region shell partition:

\begin{enumerate}[itemsep=1pt,topsep=2pt]
  \item \textbf{Deep interior}: every cell sufficiently far from
        every physical face uses fixed \texttt{D4CEN} on all axes.
        No branches in the kernel body --- this is the hot path.
  \item \textbf{Shell A}: $z$-shell rows ($k$ in the SBP closure
        zone), full $(j, i)$.  Runtime stencil dispatch on all
        three axes via \texttt{stencil\_at} +
        \texttt{apply\_stencil}.  Per-face SAT penalties added
        when the cell sits on a physical boundary row.
  \item \textbf{Shell B}: $y$-shell rows, $z$ in deep interior,
        full $i$.  $z$-stencil known to be \texttt{D4CEN}; $x, y$
        dispatched at runtime.
  \item \textbf{Shell C}: $x$-shell rows, $y, z$ in deep interior.
        Only $x$ dispatched at runtime.
\end{enumerate}
The three shell regions partition the boundary shell exactly
once, so each cell is visited at most twice (once interior, once
shell) and each SAT face fires at most once per RK4 stage at any
given cell, even at edges and corners.  The same partition appears
in \texttt{maxwell\_constraints} for the diagnostic divergences
$c_D, c_B$.

\subsection{Ghost-zone exchange}

\texttt{sync\_vars(NGFS\&, var\_type, kidx)} performs the standard
non-blocking neighbour exchange in three sequential axis sweeps
($x$, then $y$, then $z$).  Pack and unpack are
\texttt{parallel\_for} kernels writing into a
\texttt{Kokkos::View<double*>} per axis face; on the CUDA backend
the device buffer is deep-copied to a host mirror before the MPI
call (or sent directly when built with
\texttt{MAXWELL\_CUDA\_AWARE\_MPI=ON}).

The axis order is sequential rather than fully parallel because
the $y$-axis send box includes $x$-ghost columns at $j \in [g_s,
2g_s)$; the $y$-sync only delivers correct corner data if $x$ has
already been synced.  The Maxwell production stencils are
axis-aligned and never read corner ghosts, so a fully parallel
exchange would be functionally correct for the PDE solve, but the
\texttt{test\_sync} harness corrupts and verifies all ghost cells
including corners.  Sequential order matches the C version.

The special form \texttt{sync\_vars(gfs, EVOLVED, -1)} syncs the
\texttt{state} buffer instead of \texttt{K[kidx]} --- used by the
startup sync test in the driver to validate
\texttt{set\_initial\_data}.

\subsection{Reductions}

\texttt{l2\_error\_analytic} is one
\texttt{Kokkos::parallel\_reduce} accumulating the squared $L^2$
residual across all nine evolved fields into a single
\texttt{double}, followed by \texttt{MPI\_Allreduce} on
$(\text{sum}, \text{count})$.  The collective is the only
synchronisation between ranks during evaluation.

\subsection{File map}

\begin{center}
\small
\begin{tabular}{@{}lll@{}}
\toprule
File & Origin & Role \\
\midrule
\texttt{driver.cpp}            & ported    & Top-level driver, RK4 loop \\
\texttt{maxwell\_eqs.\{hpp,cpp\}} & ported & Time deriv + constraints + IC + L2 \\
\texttt{rk4.\{hpp,cpp\}}       & ported    & Kokkos RK4 \\
\texttt{numerical.\{hpp,cpp\}} & ported    & \texttt{apply\_dissipation} \\
\texttt{comm.\{hpp,cpp\}}      & ported    & Ghost-zone sync \\
\texttt{gf.\{hpp,cpp\}}        & ported    & View-owning \texttt{NGFS} \\
\texttt{io.\{hpp,cpp\}}        & ported    & HDF5 output + checkpoint/restart \\
\texttt{HDF5BinaryWrite.\{hpp,cpp\}} & ported & HDF5 file/group/dataset wrapper \\
\texttt{analytic\_parameters.h}  & ported    & POD parameter structs (per-source) \\
\texttt{analytic\_solutions.hpp} & ported    & Kokkos-callable analytic evaluators \\
\texttt{derivatives.hpp}         & ported    & D4CEN, SBP-4-2 closures, dissipation \\
\texttt{simple\_maxwell.h}     & verbatim  & SymPy-generated SAT \& interior macros \\
\texttt{generate\_ccode.py}    & verbatim  & SymPy emitter \\
\texttt{domain.\{c,h\}}        & verbatim  & MPI Cartesian topology \\
\texttt{maxwell\_parameters.\{h,cpp\}} & verbatim & POD parameter ABI;
                                                    TOML $\to$ \texttt{maxwell\_param\_st} \\
\texttt{parameter.\{hpp,cpp\}}  & verbatim  & Internal toml++ helpers
                                              (namespace \texttt{parameters}) \\
\texttt{toml.hpp}              & verbatim  & toml++ vendored header \\
\texttt{timer.\{c,h\}}         & verbatim  & Hierarchical timer \\
\texttt{tests/test\_topology.c}& verbatim  & MPI decomposition unit test \\
\texttt{tests/test\_rk4.cpp}   & ported    & RK4 4th-order convergence test \\
\texttt{tests/test\_sync.cpp}  & ported    & Ghost-zone exchange unit test \\
\bottomrule
\end{tabular}
\end{center}

\section{Build}
\label{sec:build}

\subsection{Prerequisites}

\begin{itemize}[itemsep=2pt,topsep=2pt]
  \item CMake $\geq 3.18$.
  \item C++20 compiler.  Tested: GCC 13, GCC 15.  GCC 12 has a
        known interaction with toml++; see
        \S\ref{sec:known}.
  \item MPI (OpenMPI, MPICH, or any MPI-3 implementation).
  \item HDF5 with the C and HL components.
  \item For the GPU build: CUDA Toolkit $\geq 11.4$ targeting an
        appropriate GPU architecture.
\end{itemize}

\subsection{Default OpenMP build}
\begin{lstlisting}[style=shell]
cd Maxwell_Kokkos/src
cmake -S . -B build
cmake --build build -j
\end{lstlisting}
This produces \texttt{build/maxwell\_system} plus the three test
executables under \texttt{build/tests/}.

\subsection{CUDA build}
\begin{lstlisting}[style=shell]
cmake -S . -B build \
      -DKokkos_ENABLE_OPENMP=OFF \
      -DKokkos_ENABLE_CUDA=ON   \
      -DKokkos_ARCH_AMPERE80=ON
cmake --build build -j
\end{lstlisting}
Replace \texttt{Kokkos\_ARCH\_AMPERE80} with the
\texttt{Kokkos\_ARCH\_*} variable for your GPU's compute capability
(\texttt{HOPPER90} for an H100, \texttt{TURING75} for a T4, etc.).

\subsection{CUDA-aware MPI}
\begin{lstlisting}[style=shell]
cmake -S . -B build \
      -DKokkos_ENABLE_OPENMP=OFF \
      -DKokkos_ENABLE_CUDA=ON   \
      -DKokkos_ARCH_AMPERE80=ON \
      -DMAXWELL_CUDA_AWARE_MPI=ON
\end{lstlisting}
Skips the device--host buffer hop in \texttt{comm.cpp} and hands
device pointers directly to MPI.  Requires an MPI install built
with CUDA awareness (e.g.\ OpenMPI \texttt{--with-cuda}).

\subsection{CMake options reference}

\begin{center}
\small
\begin{tabular}{@{}lll@{}}
\toprule
Option & Default & Effect \\
\midrule
\texttt{Kokkos\_ENABLE\_OPENMP}    & \texttt{ON}      & OpenMP host backend \\
\texttt{Kokkos\_ENABLE\_CUDA}      & \texttt{OFF}     & CUDA device backend \\
\texttt{Kokkos\_ARCH\_<arch>}      & unset            & Required when CUDA \\
\texttt{MAXWELL\_CUDA\_AWARE\_MPI} & \texttt{OFF}     & Device pointers to MPI \\
\texttt{CMAKE\_BUILD\_TYPE}        & \texttt{Release} & Standard CMake \\
\texttt{CMAKE\_CXX\_STANDARD}      & \texttt{20}      & Hard requirement \\
\bottomrule
\end{tabular}
\end{center}

\subsection{Source-level switches}

Defining \texttt{TIMING\_ONLY} on the compile line skips HDF5
output and L$^2$ reporting in the main loop, leaving only kernel
work for clean per-step timing.

\section{Running}
\label{sec:running}

\begin{lstlisting}[style=shell]
mpirun -np <N> ./maxwell_system path/to/maxwell.toml
\end{lstlisting}

\subsection{Parameter file structure}
\label{sec:toml}

\paragraph{[grid]}
\begin{lstlisting}[style=shell]
nx = 64          # global grid points in x (total, not per rank)
ny = 64
nz = 128
x0 = -1.0        # lower corner
y0 = -1.0
z0 =  0.0
xn =  1.0        # upper corner
yn =  1.0
zn =  4.0
periodic_x = false  # if true: no SAT on this axis, ghost zones wrap
periodic_y = false
periodic_z = false
\end{lstlisting}

\paragraph{[solver]}
\begin{lstlisting}[style=shell]
ghost_size       = 3
cfl_factor       = 0.2    # <= 0.2 for all-physical; <= 0.5 for periodic x,y
max_iterations   = 2049
output_every     = 32
checkpoint_every = 512
max_checkpoints  = 2
recover          = false
use_dissipation  = false  # Kreiss-Oliger; off by default
diss_coeff       = 0.1
tau              = 1.0    # SAT penalty strength, >= 1/2
\end{lstlisting}

\paragraph{[physics]}
\begin{lstlisting}[style=shell]
kappa_D = 0.1
kappa_B = 0.1
\end{lstlisting}

\paragraph{[source]}
\begin{lstlisting}[style=shell]
[source]
type = "gaussian_beam"   # or "plane_wave" or "te_waveguide_mode"

[source.plane_wave]
ax     = 1.0       # x-polarisation amplitude
ay     = 1.0       # y-polarisation amplitude
k      = 4.0       # wavenumber
bump_a = 0.0       # envelope turn-on start (in t-z)
bump_b = 2.0       # envelope turn-on end

[source.gaussian_beam]
w0        = 0.20   # waist radius
z_waist   = 2.0    # z location of waist (inside domain for lens demo)
k         = 15.0   # wavenumber
amplitude = 1.0    # peak field amplitude
ramp_a    = 0.0    # turn-on start (in t)
ramp_b    = 1.0    # turn-on end

[source.te_waveguide_mode]
l = 2
m = 2
n = 2
\end{lstlisting}

\paragraph{[material]}
\begin{lstlisting}[style=shell]
[material]
epsilon_type = "elliptical"   # or "uniform"

[material.background]
epsilon = 1.0                 # baseline (asymptote of elliptical profile)
mu      = 1.0
sigma   = 0.0

[material.elliptical]
max = 4.0                     # peak permittivity at lens centre
x0  = 0.0
z0  = 2.0
s   = 2.0                     # steepness
a   = 0.5                     # semi-axis x
b   = 0.25                    # semi-axis z
\end{lstlisting}

\subsection{Two canonical configurations}

\paragraph{Option B: plane-wave convergence reference.}
Periodic $x, y$, non-periodic $z$, plane-wave source.  The only
configuration in which an exact, non-dispersive travelling-wave
solution exists for the plane-wave source:
\begin{lstlisting}[style=shell]
[grid]
periodic_x = true
periodic_y = true
periodic_z = false

[source]
type = "plane_wave"
\end{lstlisting}
On a sequence of refined grids, the $L^2$ error against the
analytic reduces at the order of the SBP-$(4,2)$ boundary closure
(formally 2nd-order locally at the boundary, integrating to
higher order globally; empirically around 3rd-order).

\paragraph{Option A: beam through lens.}
All six faces non-periodic with SAT, Gaussian-beam source,
elliptical permittivity:
\begin{lstlisting}[style=shell]
[grid]
periodic_x = false
periodic_y = false
periodic_z = false

[source]
type = "gaussian_beam"

[material]
epsilon_type = "elliptical"
\end{lstlisting}

\subsection{Output layout}

Each rank writes a self-describing \texttt{3D\_rank\_<R>.h5}
containing the full state at every output interval, grouped by
output index.  Groups $/0, /1, /2, \ldots$ correspond to the
initial state, the first output step, etc.\ (step $N$ is at time
$N \cdot \mathtt{output\_every}\cdot dt$).  Each group has
datasets for the nine evolved variables and the five auxiliaries
(including the diagnostic divergences $c_D, c_B$).  The
\texttt{/metadata} group carries per-rank attributes describing
the grid (origin, spacing, local and global sizes, ghost width,
variable names) --- everything needed to reassemble the global
domain.

With the beam source's smooth turn-on, group \texttt{/0} contains
identically zero fields.  The first non-trivial output is at
$t = \mathtt{output\_every}\cdot dt$; by $t \geq \mathtt{ramp\_b}$
the envelope has saturated.

An $L^2$ diagnostic \texttt{l2\_norm.dat} is written by rank 0 at
every output step, containing
$(t,\ \|U_{\text{numerical}} - U_{\text{analytic}}\|_{L^2})$.
For \texttt{plane\_wave} with periodic $x, y$ and for
\texttt{te\_waveguide\_mode}, this is a real convergence metric.
For \texttt{gaussian\_beam} the analytic reference is only the
paraxial approximation, so the $L^2$ error reaches a paraxial-scale
plateau, not zero.

\subsection{Visualization}

\begin{lstlisting}[style=shell]
python scripts/make_xdmf.py --dir <run_dir> --dt-per-output T
\end{lstlisting}
where
$T = \mathtt{output\_every}\cdot c\cdot \min(dx, dy, dz)$ is the
physical time between output groups.  The script scans the
per-rank HDF5 files, reads their \texttt{/metadata}, and emits a
single \texttt{maxwell.xmf} alongside.  The manifest is pure XML
metadata --- the payload bytes stay in HDF5.  The result opens
directly in:
\begin{itemize}[itemsep=1pt,topsep=2pt]
  \item ParaView (Xdmf3 reader preferred),
  \item VisIt (Xdmf plugin),
  \item PyVista (\texttt{pyvista.XdmfReader("maxwell.xmf")}).
\end{itemize}
A time slider stepping through the $N$ output steps appears
automatically; all 14 fields are selectable as point-data
attributes.

\paragraph{Grid-layout details.}
The XDMF uses a \texttt{3DRectMesh} topology with explicit
\texttt{VXVYVZ} geometry (one per-axis coordinate list).  This
avoids the long-standing XDMF \texttt{ORIGIN\_DXDYDZ} ambiguity in
which VisIt and ParaView disagree on whether the origin tuple is
in spatial $(x, y, z)$ or topology $(z, y, x)$ order;
\texttt{VXVYVZ} tags each axis explicitly.

\paragraph{Legacy fallback.}
\texttt{scripts/hdf5\_to\_vtk.py} reassembles the global grid on
rank 0 and writes one VTI per time step.  Serial, duplicates the
data on disk; supplied only as a fallback.

\subsection{Checkpoint restart}

Set \texttt{[solver].recover = true} and re-run with the same
TOML.  The driver locates the highest-iteration checkpoint for
this rank and resumes; HDF5 group numbering continues monotonically
from the recovered iteration.

\subsection{Diagnostics and common problems}

\paragraph{$L^2$ error blows up exponentially.}
Usually CFL is too large.  Reduce \texttt{cfl\_factor} to 0.2 (six-
face SAT) or verify the configuration is in fact Option B if you
expected $\mathrm{cfl} = 0.4$ to work.

\paragraph{$D$ field looks empty in visualization.}
Check the time slice.  Group \texttt{/0} is the initial state,
which is zero for both plane-wave and Gaussian-beam sources.  For
the default ramp $[0, 1]$ on the shipped TOML, the beam reaches
${\sim}3\%$ by group \texttt{/1} ($t = 0.2$) and full amplitude by
$t \geq \mathtt{ramp\_b}$; visualise group \texttt{/5} or later
for the default parameters.

\paragraph{Paraxial residual looks larger than expected.}
Verify $k\cdot w_0 \gtrsim 5$.  Residual scales as $1/(kw_0)^2$.

\paragraph{Beam appears to reflect off the transverse walls.}
The $x, y$ faces use $g = 0$ SAT; reflection amplitude is
proportional to the field amplitude at the wall.  Increase $L_x,
L_y$ or reduce $w_0$ so the tail is negligible.  Rule of thumb:
$L_x \geq 10\, w_0$, where $w_0$ is the local spot size along the
propagation, not the waist.

\paragraph{High-frequency oblique wavefronts at late times.}
Combined signature of (a) a real standing wave from incident +
lens-reflected beam interfering every $\lambda/2 = \pi/k$ and (b)
spurious angular reflections of lens-scattered energy off every
SAT face.  (a) is physics; (b) is the SAT limitation
(\S\ref{sec:oblique-reflection}).  Enabling Kreiss-Oliger
dissipation suppresses the high-frequency component.  For truly
angle-independent absorption the code would need a PML
(\S\ref{sec:pml}).

\section{Tests}
\label{sec:tests}

Three test executables ship with the build, plus a Python
convergence harness (\S\ref{sec:convergence}).

\subsection{\texttt{test\_topology}}

Pure CPU; exercises \texttt{automatic\_topology}, the greedy
prime-factor MPI Cartesian decomposition.  Single-process; runs
through a battery of grid shapes and process counts.

\subsection{\texttt{test\_rk4}}

Solves $dy/dt = -y$ with $y(0) = 1$ on a small grid for
$t \in [0, 1]$ at two time-step sizes.  Asserts the error ratio
between $\Delta t = 0.01$ and $\Delta t/2$ is between 12 and 20,
i.e.\ within tolerance of the expected fourth-order rate
$2^4 = 16$.  Drives the Kokkos RK4 + sync paths end-to-end on a
trivial PDE.

\subsection{\texttt{test\_sync}}

The infrastructure correctness test.  Workflow: fill all evolved
and auxiliary fields with a known smooth function on the local
grid, corrupt every ghost cell with a sentinel value, call
\texttt{sync\_vars}, then verify every cell (including ghosts)
matches the analytic value.  Driven by \texttt{run\_tests.sh} at
1, 2, 4, 6, 8, and 27 ranks across three boundary configurations:
fully physical (\texttt{0 0 0}), periodic $x, y$ + physical $z$
(\texttt{1 1 0}, the Maxwell default), and fully periodic
(\texttt{1 1 1}).

\subsection{Running the suite}

\begin{lstlisting}[style=shell]
cd build
ctest                  # runs tests/run_tests.sh
                       # = test_topology, test_rk4, and the
                       #   12-config test_sync grid

cd ..
python scripts/convergence_test.py \
       --solver src/build/maxwell_system \
       --resolutions 16 24 32 48 \
       --final-time 0.1 --np 4
\end{lstlisting}

\section{Verification: convergence testing}
\label{sec:convergence}

A direct test for the correctness of the interior stencil, each
face's SBP-SAT closure, and the additive penalty at edges and
corners is obtained by forcing the numerical solution to converge
to a known analytic solution and measuring the convergence rate.

\subsection{Test problem}

The test uses the \texttt{te\_waveguide\_mode} (\S\ref{sec:te-mode}),
which is an exact vacuum Maxwell solution with
$\eps = \mu = 1, \sigma = 0$ on the unit cube $[0,1]^3$.  When
\texttt{[source].type = "te\_waveguide\_mode"}, the code routes
the analytic as boundary data $g$ on \emph{every} physical face
(not just lower-$z$).  At each SAT-face point, the scheme
evaluates the analytic at the local $(x, y, z, t)$ and passes it
into the relevant \texttt{APPLY\_SAT\_*} macro.  The numerical
state is driven toward the analytic through the
incoming-characteristic projection of every physical face; the
outgoing characteristics are untouched, so the test is consistent
with the method's energy argument.

The scheme being tested is:
\begin{itemize}[itemsep=1pt,topsep=2pt]
  \item Deep interior: 4th-order centered stencil \texttt{D4CEN}
        in all three directions.
  \item Near-boundary rows on physical axes: SBP 4-2 one-sided
        closures.
  \item Exactly one boundary row per physical face: SAT penalty
        with $g$ from the analytic.
  \item Edges (two faces meet): two SAT penalties additive.
  \item Corners (three faces meet): three SAT penalties additive.
  \item Time integration: classical RK4.
\end{itemize}

\subsection{The four configurations}

The boundary periodicities are varied per axis to change which
boundary structures are exercised:

\begin{center}
\small
\begin{tabular}{lccccc}
\toprule
Configuration       & \texttt{periodic\_x} & \texttt{periodic\_y} & \texttt{periodic\_z}
                    & Faces / Edges / Corners & Expected $L^2$ rate \\
\midrule
\texttt{fully\_periodic} & true  & true  & true  & 0 / 0 / 0   & 4 \\
\texttt{z\_physical}     & true  & true  & false & 2 / 0 / 0   & 3 \\
\texttt{yz\_physical}    & true  & false & false & 4 / 4 / 0   & 3 \\
\texttt{all\_physical}   & false & false & false & 6 / 12 / 8  & 3 \\
\bottomrule
\end{tabular}
\end{center}

\paragraph{What each configuration localises.}
\begin{itemize}[itemsep=2pt,topsep=2pt]
  \item \texttt{fully\_periodic}: no SBP closures, no SAT.  All
        derivatives are \texttt{D4CEN} with periodic ghosts.
        Tests only the interior stencil.  Expected rate 4.
  \item \texttt{z\_physical}: only the two $z$-faces are physical;
        $x, y$ stay periodic.  Tests interior + $z$-face SBP +
        SAT.  No physical edges or corners.  A drop here means
        face-level SBP or SAT is broken.
  \item \texttt{yz\_physical}: adds four $y$-$z$ edges.  A drop
        from \texttt{z\_physical} to \texttt{yz\_physical} means
        the additive edge penalty is broken.
  \item \texttt{all\_physical}: full setup --- six faces, twelve
        edges, eight corners.  A drop here means the three-face
        corner penalty is broken.
\end{itemize}

\subsection{Why the expected rates are 4 and 3}

The scheme is formally:
\begin{itemize}[itemsep=1pt,topsep=2pt]
  \item 4th-order in the interior (\texttt{D4CEN} is $O(h^4)$).
  \item 2nd-order locally at the SBP-4-2 boundary rows (Strand
        closures are $O(h^2)$).
\end{itemize}
For a hyperbolic system discretized by an SBP operator with
interior order $2p$ and boundary truncation order $s$, the
classical Svärd--Nordström result predicts a global $L^2$ rate of
$\min(2p,\, s + 1)$ at late times.  For SBP-4-2 with $2p = 4,
s = 2$: rate 3.  The interior-only configuration sidesteps the
boundary closures and retains the full 4th-order rate.

Two practical notes:
\begin{itemize}[itemsep=2pt,topsep=2pt]
  \item The $s + 1$ rule applies per face; edges and corners do
        not further degrade the rate because the additive SAT
        gives independent energy estimates that combine as a
        maximum, not a product.  This is the feature being tested.
  \item Measured rates often approach 3 from above at finite $N$
        --- rates of 3.1--3.2 at $N = 48$ are typical for the
        SBP-4-2 scheme; they asymptote to exactly 3 as
        $N \to \infty$.  A rate persistently \emph{below} 3 at
        moderate $N$ indicates a real problem.
\end{itemize}

\subsection{Running the convergence test}

\begin{lstlisting}[style=shell]
python scripts/convergence_test.py \
    --solver src/build/maxwell_system \
    --resolutions 16 24 32 48 \
    --final-time 0.1 \
    --np 4
\end{lstlisting}
Options: \texttt{--solver PATH}, \texttt{--resolutions N1 N2 ...},
\texttt{--final-time T}, \texttt{--cfl C} (default 0.2),
\texttt{--np N}, \texttt{--work-dir DIR}.

The driver runs the four configurations in sequence; each prints
per-resolution errors and pairwise rates.  A summary table at the
end compares measured final rates against expected:

\begin{verbatim}
=== SUMMARY ===
config              N=16        N=24        N=32        N=48  final rate
fully_periodic  0.001041   0.0002078   6.597e-05   1.306e-05    3.994   OK
z_physical       0.00739    0.002273   0.0009391    0.000261    3.158   OK
yz_physical     0.009489    0.002867    0.001174    0.000325    3.167   OK
all_physical     0.01092    0.003261    0.001328    0.000368    3.166   OK
\end{verbatim}

A rate is flagged \texttt{SUSPICIOUS} if it deviates from expected
by more than 1.0; the driver exits nonzero.

\subsection{Baseline results (Kokkos OpenMP backend, validated)}

\begin{center}
\small
\begin{tabular}{lcccc|c}
\toprule
Configuration            & $N=16$               & $N=24$              & $N=32$             & $N=48$              & Final rate \\
\midrule
\texttt{fully\_periodic} & $1.04\cdot 10^{-3}$  & $2.08\cdot 10^{-4}$ & $6.60\cdot 10^{-5}$ & $1.31\cdot 10^{-5}$ & 3.99 \\
\texttt{z\_physical}     & $7.39\cdot 10^{-3}$  & $2.27\cdot 10^{-3}$ & $9.39\cdot 10^{-4}$ & $2.61\cdot 10^{-4}$ & 3.16 \\
\texttt{yz\_physical}    & $9.49\cdot 10^{-3}$  & $2.87\cdot 10^{-3}$ & $1.17\cdot 10^{-3}$ & $3.25\cdot 10^{-4}$ & 3.17 \\
\texttt{all\_physical}   & $1.09\cdot 10^{-2}$  & $3.26\cdot 10^{-3}$ & $1.33\cdot 10^{-3}$ & $3.68\cdot 10^{-4}$ & 3.17 \\
\bottomrule
\end{tabular}
\end{center}

These match the C-version's reference numbers from
\texttt{Maxwell\_Penalty/Review.md \S 1} to four significant
figures.  Run parameters: $T_{\text{final}} = 0.1$,
$\mathrm{cfl} = 0.2$, $\tau = 1.0$, four MPI ranks, uniform
$\eps = \mu = 1$, $\sigma = 0$, mode $(l, m, n) = (2, 2, 2)$.

\subsection{Interpreting a failed test}

If a future change causes a rate to drop, the configuration
hierarchy localises the fault:

\begin{itemize}[itemsep=2pt,topsep=2pt]
  \item \texttt{fully\_periodic} drops below 4: the interior
        \texttt{D4CEN} stencil, ghost-zone sync, or periodic
        wrap is broken.
  \item \texttt{z\_physical} drops below 3 while
        \texttt{fully\_periodic} stays at 4: the $z$-face SBP
        closures or the lower/upper-$z$ SAT penalty is broken.
  \item \texttt{yz\_physical} drops below 3 while
        \texttt{z\_physical} is fine: the additive edge penalty
        at $y$-$z$ edges is broken.
  \item \texttt{all\_physical} drops below 3 while
        \texttt{yz\_physical} is fine: the three-face corner
        penalty is broken.
\end{itemize}

A successful pass on all four rows confirms that every part of
the algorithm --- interior stencil, face closure, face SAT, edge
additivity, corner additivity --- is correctly implemented.

\paragraph{A subtle point: why the oblique-reflection failure mode
of \S\ref{sec:oblique-reflection} does not break this test.}
The \texttt{te\_waveguide\_mode} has nontrivial structure in every
direction and hits the $\pm x, \pm y$ faces at oblique angles ---
the exact configuration mis-handled by the $g = 0$ SAT used in the
lens demo.  One could reasonably expect spurious reflections to
pollute the rate.

This does not happen because the convergence test uses
$g = \text{analytic}$ on every physical face, not $g = 0$.  The
absorbing-boundary failure mode arises only when the prescribed
incoming data disagrees with the true physical solution at the
boundary: with $g = 0$ and a nonzero tangential field, the SAT
pulls the wave toward zero instead of toward its true value,
generating the spurious reflection.  When $g$ equals the exact
analytic at every boundary point, the SAT imposes the correct
incoming data on every face regardless of incidence angle.
Mathematically, the SAT penalty is a weak-form Dirichlet
imposition of the incoming characteristic; with the correct
right-hand side it has no angular blind spot.

The convergence test therefore verifies that the SBP-SAT scheme,
\emph{given the correct incoming data}, converges to the analytic
at the expected rate.  It is a test of the scheme's
implementation, not of the quality of absorbing boundaries in
the lens-demo sense.

\paragraph{Caveats.}
\begin{itemize}[itemsep=2pt,topsep=2pt]
  \item The \texttt{te\_waveguide\_mode} source requires
        $\eps = \mu = 1, \sigma = 0$.  The test does not
        exercise spatially-varying permittivity; that's a separate
        physics question.
  \item The constraint fields $\Psi_D, \Psi_B$ are identically
        zero in the analytic.  The constraint-damping blocks of
        the SAT are exercised only by their coupling to the EM
        blocks through residual numerical violations, so a bug
        localised entirely to those blocks might not be caught at
        full sensitivity.
  \item The test uses $g = \text{analytic}$ on every face and
        therefore does \emph{not} measure the quality of the
        $g = 0$ absorbing boundary on oblique waves --- that is a
        physics question (\S\ref{sec:oblique-reflection}) and a
        design question about replacing SAT with PML
        (\S\ref{sec:pml}), not an implementation question about
        the SAT itself.
\end{itemize}

\section{Validation status}
\label{sec:status}

\subsection{Verified on the OpenMP backend}

\begin{itemize}[itemsep=2pt,topsep=2pt]
  \item All 14 \texttt{run\_tests.sh} configurations pass:
        \texttt{test\_topology}, \texttt{test\_rk4}, and 12
        \texttt{test\_sync} cases at 1--27 ranks across three
        boundary configurations.
  \item Convergence test passes at all four configurations to
        within 5\% of the expected asymptote at $N=16, 24, 32,
        48$: 3.99 (fully periodic), 3.16/3.17/3.17 (boundary
        configurations) at $N=48$.  These rates match the
        C-version's published reference numbers in
        \texttt{Maxwell\_Penalty/Review.md \S 1} to four
        significant figures.
  \item Smoke runs of both demo TOMLs (\texttt{maxwell.toml},
        \texttt{maxwell\_waveguide.toml}) complete normally with
        the documented $L^2$ behaviour.
\end{itemize}

\subsection{Built but not yet runtime-verified}

\begin{itemize}[itemsep=2pt,topsep=2pt]
  \item \textbf{CUDA backend.}  Builds cleanly on an A100 host
        with the toolchain pitfalls in \S\ref{sec:known} resolved.
        Runtime convergence test on the GPU is the recommended
        next validation step.
  \item \textbf{CUDA-aware MPI.}  The
        \texttt{MAXWELL\_CUDA\_AWARE\_MPI=ON} compile path
        compiles, but no test distinguishes a CUDA-aware build
        from the default host-mirror build other than performance.
  \item \textbf{Checkpoint restart.}  Both
        \texttt{write\_checkpoint} and \texttt{read\_checkpoint}
        are ported and link, but no automated round-trip
        verification is in the test suite.
  \item \textbf{\texttt{make check-generated}.}
        \texttt{generate\_ccode.py} is unchanged from the C
        version and emits a \texttt{simple\_maxwell.h}
        byte-for-byte identical to the checked-in copy, but no
        CMake target exercises this.
\end{itemize}

\subsection{Recommended GPU-side test sequence}

After a successful CUDA build, validate in this order:
\begin{enumerate}[itemsep=2pt,topsep=2pt]
  \item \texttt{./tests/test\_rk4} --- expect
        \texttt{Kokkos exec space: Cuda} and the 4th-order ratio
        check to pass.
  \item \texttt{./tests/test\_sync 32 32 32 0 0 0} at
        \texttt{-np 1} --- exercises the device pack/unpack and
        the host-mirror MPI hop.
  \item \texttt{python scripts/convergence\_test.py --np 1
        --resolutions 16 24 32 --final-time 0.1} --- the strongest
        single check.  Each rate should match the OpenMP
        baseline to better than 0.5\%.
  \item Diff a 50-step \texttt{maxwell.toml} run's
        \texttt{l2\_norm.dat} between OpenMP and CUDA --- should
        agree to $10^{-12}$ or better at every output step.
\end{enumerate}

\section{Known issues and toolchain workarounds}
\label{sec:known}

\subsection{OpenMPI 4.x C++ bindings vs.\ \texttt{nvcc}}

OpenMPI's \texttt{mpi.h} pulls in \texttt{mpi/cxx/functions.h} and
\texttt{mpi/cxx/file.h}, headers from the deprecated MPI-2 C++
namespace.  Those headers declare member functions like
\texttt{MPI::Init()} inside an \texttt{extern "C"} block with
overloads --- legal under \texttt{g++} (which silently picks one),
illegal under \texttt{nvcc}'s stricter linkage check.

\textbf{Fix} (already in \texttt{CMakeLists.txt}):
\begin{lstlisting}[style=cpp]
add_compile_definitions(OMPI_SKIP_MPICXX MPICH_SKIP_MPICXX)
\end{lstlisting}
Both macros are recognised by their respective MPI implementations
and skip the entire \texttt{mpi/cxx/*.h} tree.  The code uses only
the C MPI API, so dropping the C++ bindings has no functional
effect.

\subsection{GCC 12 + toml++ + C++20}

GCC 12 (and only GCC 12) miscompiles a generic-lambda
\texttt{if constexpr} pattern used by toml++:
\begin{lstlisting}[style=cpp]
using return_type = decltype(visitor(...));
if constexpr (impl::is_constructible_or_convertible<bool, return_type>)
\end{lstlisting}
The compiler mangles \texttt{return\_type} as \texttt{\_\_T<N>} and
then fails to find the name in the \texttt{if constexpr}
expression.

\textbf{Fix (design-level).}  The original parameter header served
two unrelated roles: declaring the maxwell-specific POD parameter
ABI (structs + \texttt{parse\_maxwell\_parameters()}) and
providing toml++-typed helpers
(\texttt{namespace parameters \{ \ldots \}}) used internally by the
parser implementation.  These were split into two headers:
\begin{itemize}[itemsep=2pt,topsep=2pt]
  \item \texttt{maxwell\_parameters.h} --- public ABI: the POD
        structs the rest of the solver consumes plus the
        C-callable \texttt{parse\_maxwell\_parameters()}.  No
        toml++.  Included by \texttt{driver.cpp} (transitively
        via \texttt{maxwell\_eqs.hpp}) and any consumer of the
        parsed result.
  \item \texttt{parameter.hpp} --- internal header of the parser:
        \texttt{namespace parameters} helpers built on
        \texttt{toml::table}.  Pulls \texttt{toml.hpp}.  Included
        only by \texttt{parameter.cpp} (which implements the
        namespace) and \texttt{maxwell\_parameters.cpp} (which uses
        it to populate the structs).
\end{itemize}
toml++ exposure is now implicit in including
\texttt{parameter.hpp}; nothing else needs it.  Translation units
that never include \texttt{parameter.hpp} never see the buggy
toml++ lambda.  Side benefit: a 15{,}000-line header is no longer
pulled into the kernel-heavy translation units, which
substantially shortens compile time.

If a future patch invokes the toml++ helpers from a new
translation unit and the GCC 12 bug resurfaces, the in-place
workaround in \texttt{toml.hpp} is to bind the trait to a local
\texttt{constexpr bool}:
\begin{lstlisting}[style=cpp]
using return_type = decltype(...);
constexpr bool _returns_bool =
    impl::is_constructible_or_convertible<bool, return_type>;
if constexpr (_returns_bool) { ... }
\end{lstlisting}
Apply at both \texttt{toml.hpp:8239} and \texttt{:8258}.

\subsection{Single-GPU multi-rank}

By default each MPI rank tries to bind device 0.  CUDA contexts
can be shared, so this is correct, but performance is degraded
and some operations may serialise.  Workarounds:
\begin{itemize}[itemsep=2pt,topsep=2pt]
  \item \texttt{-np 1} for single-GPU correctness checks.
  \item NVIDIA Multi-Process Service
        (\texttt{nvidia-cuda-mps-control -d}) for $\leq 2$--$4$
        ranks.
  \item On multi-GPU nodes, one rank per GPU
        (\texttt{mpirun --map-by ppr:1:gpu}).
\end{itemize}

\subsection{C++20 hard requirement}

Kokkos 5.x rejects C++17 outright.  The C version's TOML files
build under C++17, but the Kokkos build needs at least C++20.
GCC 13+, Clang 16+, NVHPC 24+ all satisfy this.

\section*{References}

\begin{itemize}[itemsep=2pt,topsep=2pt]
  \item Strand, B.\ (1994), \emph{Summation by parts for finite
        difference approximations of $d/dx$}, J.\ Comput.\ Phys.\
        110, 47--67.  Source for the SBP-4-2 norm coefficients
        and one-sided closures used in \texttt{derivatives.hpp}.
  \item Carpenter, M., Gottlieb, D., Abarbanel, S.\ (1994),
        \emph{Time-stable boundary conditions for finite-
        difference schemes solving hyperbolic systems: methodology
        and application to high-order compact schemes},
        J.\ Comput.\ Phys.\ 111, 220--236.  Original SAT
        derivation.
  \item Svärd, M., Nordström, J.\ (2014), \emph{Review of
        summation-by-parts schemes for initial-boundary-value
        problems}, J.\ Comput.\ Phys.\ 268, 17--38.  Source of the
        $\min(2p, s+1)$ global $L^2$ rate estimate used in
        \S\ref{sec:convergence}.
  \item Munz, C.-D., Omnes, P., Schneider, R., Sonnendrücker, E.,
        Voss, U.\ (2000), \emph{Divergence correction techniques
        for Maxwell solvers based on a hyperbolic model},
        J.\ Comput.\ Phys.\ 161, 484--511.  The constraint-damping
        extension used here.
  \item Bérenger, J.-P.\ (1994), \emph{A perfectly matched layer
        for the absorption of electromagnetic waves},
        J.\ Comput.\ Phys.\ 114, 185--200.  Original split-field
        PML.
  \item Roden, J.\,A., Gedney, S.\,D.\ (2000), \emph{Convolutional
        PML (CPML): an efficient FDTD implementation of the
        CFS-PML for arbitrary media},
        Microwave Opt.\ Technol.\ Lett.\ 27, 334--339.  The CPML
        variant recommended for the proposed extension in
        \S\ref{sec:pml}.
  \item Kokkos team (2024). \emph{Kokkos C++ Performance
        Portability Programming EcoSystem},
        \url{https://kokkos.org}.
  \item Trott, C.\ et al.\ (2022), \emph{Kokkos 3: Programming
        Model Extensions for the Exascale Era}, IEEE TPDS.
\end{itemize}

\appendix

\section{Summary of SAT formulas}
\label{app:sat-formulas}

For reference, the full set of SAT contributions applied at the
boundary rows, derived in \S\ref{sec:sat}.  Let
$s = \tau \cdot (H^{-1}_0)/h$ be the per-axis penalty scale
(with $H^{-1}_0 = 48/17$ for the $(4,2)$ SBP norm), and define
$c, PP$ as in \S\ref{sec:characteristics}.  Where $g$ is zero, the
$\delta$ quantities reduce to the evolved fields themselves.

\paragraph{Lower-$z$ (nonzero $g$, the incoming beam):}
\begin{align*}
  \dot B_x   &\mathrel{-}= s \cdot \tfrac{c}{2} (\delta B_x - PP \cdot \delta D_y) \\
  \dot D_y   &\mathrel{-}= s \cdot \tfrac{c}{2} (-\delta B_x / PP + \delta D_y) \\
  \dot B_y   &\mathrel{-}= s \cdot \tfrac{c}{2} (\delta B_y + PP \cdot \delta D_x) \\
  \dot D_x   &\mathrel{-}= s \cdot \tfrac{c}{2} (\delta B_y / PP + \delta D_x) \\
  \dot B_z   &\mathrel{-}= s \cdot \tfrac{1}{2} (\delta B_z - \delta \Psi_B) \\
  \dot \Psi_B &\mathrel{-}= s \cdot \tfrac{1}{2} (-\delta B_z + \delta \Psi_B) \\
  \dot D_z   &\mathrel{-}= s \cdot \tfrac{1}{2} (\delta D_z - \delta \Psi_D) \\
  \dot \Psi_D &\mathrel{-}= s \cdot \tfrac{1}{2} (-\delta D_z + \delta \Psi_D)
\end{align*}

\paragraph{Upper-$z$ ($g = 0$):}
\begin{align*}
  \dot B_x   &\mathrel{-}= s \cdot \tfrac{c}{2} (B_x + PP \cdot D_y) \\
  \dot D_y   &\mathrel{-}= s \cdot \tfrac{c}{2} (B_x / PP + D_y) \\
  \dot B_y   &\mathrel{-}= s \cdot \tfrac{c}{2} (B_y - PP \cdot D_x) \\
  \dot D_x   &\mathrel{-}= s \cdot \tfrac{c}{2} (-B_y / PP + D_x) \\
  \dot B_z   &\mathrel{-}= s \cdot \tfrac{1}{2} (B_z + \Psi_B) \\
  \dot \Psi_B &\mathrel{-}= s \cdot \tfrac{1}{2} (B_z + \Psi_B) \\
  \dot D_z   &\mathrel{-}= s \cdot \tfrac{1}{2} (D_z + \Psi_D) \\
  \dot \Psi_D &\mathrel{-}= s \cdot \tfrac{1}{2} (D_z + \Psi_D)
\end{align*}

\paragraph{Lower-$x$, lower-$y$, upper-$x$, upper-$y$:} obtained
from the lower-$z$ and upper-$z$ formulas above by cyclic
permutation of the axes.  The explicit forms live in
\texttt{simple\_maxwell.h}, macros
\texttt{APPLY\_SAT\_\{LOWER,UPPER\}\_\{X,Y,Z\}}.

\end{document}
